\section{问题模型}\label{sec:formulation}
输入为微操作 DAG $G=(\mathcal V,\mathcal E)$。操作节点携带流水线、周期数与缓冲区集合；缓存管理节点为每个缓冲区 $b$ 给出唯一的分配节点 $a_b$、释放节点 $f_b$、大小 $s_b$ 与缓存类型 $\tau(b)$。基础调度是全部原始节点的一个拓扑排列。评测使用五个抽象容量：L1${}=4096$、UB${}=1024$、L0A${}=256$、L0B${}=256$、L0C${}=512$。

\begin{figure}[t]\centering
\resizebox{.98\linewidth}{!}{\begin{tikzpicture}[
      font=\footnotesize,
      op/.style={draw=ksTikzOp!85!black, fill=ksTikzOp!13, rounded corners=2.5pt,
                 minimum height=8.6mm, minimum width=21mm, align=center,
                 line width=0.7pt},
      mgmt/.style={draw=ksTikzMg!70!black, fill=ksTikzMg!9, rounded corners=2.5pt,
                   minimum height=8.6mm, minimum width=21mm, align=center,
                   line width=0.7pt},
      dep/.style={-{Stealth[length=2.4mm,width=2mm]}, line width=0.7pt,
                  black!60},
      life/.style={{Bar[width=1.5mm]}-{Bar[width=1.5mm]}, line width=1.1pt},
      note/.style={font=\footnotesize, anchor=west, text=black}
    ]
    \definecolor{ksTikzOp}{HTML}{4C72B0}
    \definecolor{ksTikzMg}{HTML}{C44E52}
    \definecolor{ksTikzGen}{HTML}{DD8452}

    % DAG nodes
    \node[mgmt] (a0) at (0,0)    {\textsc{alloc} $b_0$\\[-1.5pt] {\scriptsize\color{black}backed}};
    \node[op]   (ci) at (3.1,0)  {\textsc{copy-in}\\[-1.5pt] $\{b_0\}$};
    \node[op]   (cp) at (6.2,0)  {\textsc{compute}\\[-1.5pt] $\{b_0,b_1\}$};
    \node[op]   (co) at (9.3,0)  {\textsc{copy-out}\\[-1.5pt] $\{b_1\}$};
    \node[mgmt] (a1) at (3.1,-1.8) {\textsc{alloc} $b_1$\\[-1.5pt] {\scriptsize\color{black}generated}};
    \node[mgmt] (f0) at (6.2,-1.8) {\textsc{free} $b_0$};
    \node[mgmt] (f1) at (9.3,-1.8) {\textsc{free} $b_1$};

    % dependency edges
    \draw[dep] (a0) -- (ci);
    \draw[dep] (ci) -- (cp);
    \draw[dep] (cp) -- (co);
    \draw[dep] (a1) -- (cp);
    \draw[dep] (cp) -- (f0);
    \draw[dep] (co) -- (f1);

    % logical lifetimes with spill-class costs
    \draw[life, ksTikzOp]  (0,-3.0)  -- (6.2,-3.0);
    \node[note, anchor=north, align=center] at (3.1,-3.12)
      {$b_0$: backed ($\kappa_{b_0}{=}1$) $\Rightarrow$ $1\times s_{b_0}$ (reload only)};
    \draw[life, ksTikzGen] (3.1,-3.72) -- (9.3,-3.72);
    \node[note, anchor=north, align=center] at (6.2,-3.84)
      {$b_1$: generated ($\kappa_{b_1}{=}2$) $\Rightarrow$ $2\times s_{b_1}$ (write-back $+$ reload)};

    % legend
    \node[mgmt, minimum width=4.5mm, minimum height=3.2mm] at (0.1,-4.60) {};
    \node[note] at (0.4,-4.60) {cache management};
    \node[op, minimum width=4.5mm, minimum height=3.2mm] at (3.15,-4.60) {};
    \node[note] at (3.45,-4.60) {operation};
    \node[note, anchor=west] at (5.25,-4.60)
      {Canonical P2 identity: $\mathrm{Tr}=\mathrm{Cl}+2\mathrm{Dt}=(\mathrm{Cl}+\mathrm{Dt})+\mathrm{Dt}=\mathrm{Vol}+\mathrm{Dt}$};
  \end{tikzpicture}
}
\caption{微操作 DAG 与基准溢出类别。缓存管理节点界定每个缓冲区的逻辑生命周期：有输入备份的缓冲区溢出计费一次（仅重载），生成且无备份的缓冲区计费两次（写回加重载）。}
\label{fig:dag-example-zh}\end{figure}

评测器对缓冲区作\textbf{静态}分类：若 $b$ 出现在任意 \texttt{COPY-IN} 节点中，则 $\kappa_b=1$（有备份）；否则 $\kappa_b=2$（生成且无备份）。一次溢出事件计费 $\mathrm{tr}_b=\kappa_b s_b$。由于输入不提供逐操作的读写角色，本文不推断动态 dirty 状态；该分类是评测器的定义，而非对硬件一致性协议的断言。令两类累计溢出体积（有备份/clean、生成/dirty）为 $\mathrm{Cl},\mathrm{Dt}$，则
\begin{equation}\label{eq:extra-decomp-zh}\mathrm{Tr}=\mathrm{Cl}+2\mathrm{Dt}=\mathrm{Vol}+\mathrm{Dt},\qquad \mathrm{Vol}=\mathrm{Cl}+\mathrm{Dt},\end{equation}
故 $\mathrm{Vol}\le\mathrm{Tr}\le2\mathrm{Vol}$。该界说明固定溢出体积时，类别构成的影响至多为两倍，但拓扑序和地址规划还可改变 $\mathrm{Vol}$ 本身；紧性证明及适用边界见补充材料第 2.1 节。

物理驻留段必须占据所属缓存内的连续地址区间，且同时驻留的区间互不重叠。溢出以 \texttt{SPILL-OUT} 与 \texttt{SPILL-IN} 配对中断驻留，重载后的地址可以不同；原始使用者与基准的 \texttt{FREE} 事件均为强制驻留事件。P1 仅评价 ALLOC--FREE 的逻辑活跃峰值；P2 输出扩展调度、地址与溢出表，按键 $(\mathrm{Tr},n,T)$ 排序；P3 按 $(T,\mathrm{Tr},n)$ 排序，其中 $T$ 为流水线最早完成时间。P1 逻辑峰值不等于溢出后 P2 计划的物理峰值，P3 亦可能有意接受更多流量以换取更短时间，因此三种视图从不混合比较。所有物理结果均通过标准评测器检查节点覆盖、拓扑合法性、连续地址打包与溢出依赖；完整评测器契约及 P2--P3 配对权衡见补充材料第 1 节和第 4.4 节。
